\chapter{Conclusiones}

El desarrollo de este trabajo práctico permitió familiarizarse con el flujo de trabajo 
estándar para la programación de sistemas embebidos ARM sin depender de entornos 
integrados propietarios. Se logró comprender la importancia de los archivos de 
\textit{startup} y el script del enlazador para el inicio correcto del microcontrolador.

El uso de herramientas como Make y OpenOCD facilita la integración en sistemas de 
desarrollo flexibles y potentes, permitiendo un control total sobre el proceso de 
generación de código y depuración. La configuración directa de registros para el control 
de GPIO sentó las bases para el manejo de periféricos más complejos en futuros trabajos.
